\section*{Заключение}
\addcontentsline{toc}{section}{Заключение}

В ходе выполнения практических работ были освоены основные средства функционального программирования на языке Haskell: рекурсия, сопоставление с образцом, алгебраические типы данных, классы типов, монады, трансформеры монад и тестирование свойств.

В первой работе были реализованы базовые функции и решение задачи защиты с типом \mintinline{haskell}{Either}. Во второй работе были реализованы пользовательские типы, представители классов типов и парсеры. В третьей работе были созданы проекты для глитч-обработки изображения и поиска в лабиринте с использованием \mintinline{haskell}{RWS}. В четвёртой работе был реализован Stack-проект текстовой игры по лабиринту с собственным трансформером \mintinline{haskell}{MyStateT}. В пятой работе были реализованы функции и проверены их свойства с помощью QuickCheck.

Все основные этапы работы фиксировались в Git. Для отчёта была построена временная диаграмма коммитов и сдач, приведена структура репозитория, перечислены ключевые коммиты и отражено использование нейросетей в файлах практик.
